\documentclass[10pt,a4paper]{amsart}
\usepackage[margin=19mm]{geometry}
\usepackage{amsmath,amssymb,mathtools}
\usepackage[scr=rsfs,cal=euler]{mathalfa}
\usepackage{xfrac}
\usepackage[unicode,hidelinks,bookmarks=false]{hyperref}
\newcommand{\ie}{i.e.\ }
\newcommand{\cf}{cf.\ }
\newcommand{\resp}{respectively\ }
\newcommand{\Subsection}{\subsection}
\providecommand{\nobreakdash}{\nobreak\hbox{-}}
\setlength{\emergencystretch}{2em}
\multlinegap0pt
\usepackage{amssymb}
\usepackage{tikz}
\newtheorem{lemma}{Lemma}
\newcommand{\E}{\mathbb{E}}
\newcommand{\hF}{\tilde{F}_{p,q}}
\newcommand{\hf}{\tilde{f}_{p,q}}
\newcommand{\htc}{\hat{\tc}}
\newcommand{\n}{n_{\st}}
\newcommand{\st}{\textbf{t}}
\newcommand{\tc}{\mathcal{T}}
\title{Asymptotic normality for general subtree counts in conditioned Galton--Watson trees}
\author{Fameno Rakotoniaina, Dimbinaina Ralaivaosaona}
\date{Source: arXiv:2603.08076v1 (CC BY 4.0)}
\begin{document}
\maketitle

\noindent\emph{Source and licence.} Asymptotic normality for general subtree counts in conditioned Galton--Watson trees. Version arXiv:2603.08076v1 (CC BY 4.0), distributed by arXiv under the Creative Commons Attribution 4.0 International licence, http://creativecommons.org/licenses/by/4.0/, as recorded in the arXiv licence metadata for this exact version and shown on the abstract page https://arxiv.org/abs/2603.08076v1. Full licence notice: \texttt{licenses/arxiv-2603.08076-CC-BY-4.0.txt}.\par\smallskip

\noindent\textit{Editorial scope.} Counted unit: the article's lemma lem:toll\_add (source0.tex lines 429--435). The article's complete original proof of it is delivered with it (source0.tex lines 437--478, one \texttt{proof} environment of 42 lines). No mathematics is written, repaired or re-derived: the statement and the proof are the article's own lines, in the article's own notation, and no retained line is substituted or re-flowed.  Retained uncounted as the source-local context the proof needs: lines 155--162.  Nothing was omitted from any retained span, so the package carries no omission record.  No retained line contains a citation, so the wrapper carries no bibliography and no bibliography entry.  No retained line contains a figure environment, an image-inclusion command or drawing code, so no image is delivered and the package declares no asset dependency.  Editorial formatting only, in addition: an AMS \texttt{amsart} preamble that restates the article's own declarations and macros used by the retained lines, namely the environments \texttt{lemma} on separate counters, together with the standard packages those lines need.  The retained environments print in this document as Lemma 1 in order of appearance, whereas the article prints the same items with the numbering of its own full document.  The complete native source of the article as distributed in the arXiv e-print for this exact version is bundled unchanged as \texttt{sources/195982/source0.tex}.
\begin{lemma}\label{expectation of products of n(T) with finite tree size}
Let $\st$ and $\st'$ be fixed trees with maximum outdegree at most $\Delta$ and height at most $M$. Then the following statements hold:
\begin{itemize}
\item[(a)] If $\mathbb{E}(\xi^{\Delta+2}) < \infty$, then $\mathbb{E}(\n(\tc)\,|\tc^{(M)}|) < \infty$ and $\mathbb{E}(\n(\htc)\,|\htc^{(M)}|) < \infty$.\label{st:a}
\item[(b)] If $\mathbb{E}(\xi^{\Delta+2}) < \infty$, then $\mathbb{E}(\n(\tc_n)\,|\tc_n^{(M)}|) = O(1)$ as $n \to \infty$.
\item[(c)] If $\mathbb{E}(\xi^{2\Delta+1}) < \infty$, then $\mathbb{E}(n_{\st}(\tc_n)\,n_{\st'}(\tc_n)) = O(1)$ as $n \to \infty$.
\end{itemize}
\end{lemma}

\begin{lemma}\label{lem:toll_add}
	Assume that $\xi$ satisfies $\mathbb{E}(\xi^{2\Delta+1})<\infty$ and let $\hF$ and $\hf$ be as defined above. Then, we have
	\begin{equation}
	\E\big(\hf(\tc_n)\hF(\tc_n)\big)=O(1),
	\end{equation}
 as $n\to\infty$, where the implied constant is independent of $p$ and $q$.
\end{lemma}

\begin{proof}
We can assume that \( p \leq n < q \); otherwise, the expression inside the expectation equals zero, and there is nothing to prove.  By definition
\[
\hF(\tc_n)=\sum_{v\in\tc_n}\hf(\tc_{n,v}),
\]
where $\tc_{n,v}$ denotes the fringe subtree of $\tc_n$ rooted at $v$. The strategy is to split the sum according to the distance $d(v)$ from $v$ to the root. So, we define
	\[
		\Sigma_1 := \sum_{d(v)<M}\hf(\tc_{n,v}), \ \text{and }\  \Sigma_2 := \sum_{d(v)\geq M}\hf(\tc_{n,v}).
	\]
 
Let us look at the contribution from $\Sigma_1$ first. Observe that
\[
\E\left(\hf(\tc_n)\Sigma_1\right) = \sum_{\ell=0}^{M-1}\mathbb{E}\left(\hf(\tc_n)\sum_{d(v)=\ell}\hf(\tc_{n,v})\right).
\]
For each $\ell$, let $t_1, \ t_2, \cdots, \ t_d $ be the branches of $\st$ rooted at nodes at level $\ell$. Conditioning on $\tc_n^{(\ell)}$, let $T_1,\ T_2, \cdots ,T_k$ be the branches $\tc_n$ rooted at the nodes at level $\ell$. Further, we also condition on the sizes of $T_1,\ T_2, \cdots T_k$, say $n_1,\ n_2, \cdots , n_k$. Thus, we have
\[
\mathbb{E}\left(\hf(\tc_n)\sum_{d(v)=\ell}\hf(\tc_{n,v})\Big|\tc_n^{(\ell)}=T \wedge (\forall i \  (|T_i|=n_i)) \right)=
\]
\begin{equation}\label{expectation of centered f, proof}
    \E\left(\left(\sum_{\pi}\prod_{j=1}^{d}n_{t_j}(T_{\pi_j})-\mu_n\right)\sum_{\substack{1\leq i\leq k\\ p\leq n_i <q}}\left(\n(T_i)-\mu_{n_i}\right)\right),
\end{equation}
where the first summation inside the expectation is over possible embeddings $\pi$ of $\st^{(\ell)}$ into $T^{(\ell)}$, and for each embedding $\pi$, the node ranked $j$-th at level $\ell$ in $\st^{(\ell)}$ is mapped to the node ranked $\pi_j$ in $T^{(\ell)}$.

We know from Lemma \ref{expectation of products of n(T) with finite tree size} that $(\mu_k)_{k\geq 1}$ is a bounded sequence.  Expanding the term inside the expectation \eqref{expectation of centered f, proof} and using the fact the $T_i$'s are independent copies of conditioned Galton--Watson trees, the expectations of the non-constant terms are of the form $\E(\n(\tc_{n_i}))$,  $\E(n_{t_j}(\tc_{n_i}))$ or $\E(n_{t_j}(\tc_{n_i})\n(\tc_{n_i}))$; these are all bounded terms according to Lemma \ref{expectation of products of n(T) with finite tree size}. Hence
\[
   \mathbb{E}\left(\hf(\tc_n)\sum_{d(v)=\ell}\hf(\tc_{n,v})\Big|\tc_n^{(\ell)}=T \wedge (\forall i \  (|T_i|=n_i))\right)=O\left(w_{\ell}(T)n_{\st^{(\ell)}}(T^{(\ell)})\right)
\]
Taking the expectation again and summing over $\ell\in \{0,1,2,\cdots, M-1\}$, (noting that $w_{(\ell)}(T)\leq |T^{(\ell)}|$) we deduce from Lemma \ref{expectation of products of n(T) with finite tree size} that 
\[
\E\left(\hf(\tc_n)\Sigma_1\right)=O\left(\sum_{\ell=0}^{M-1}\E\left(n_{\st^{(\ell)}}(\tc_n)|\tc_n^{(\ell)}|\right)\right)=O(1).
\]

We take the same approach for the contribution from $\Sigma_2$, and we have  
\[
\E\left(\hf(\tc_n)\Sigma_2\right) = \sum_{\ell\geq M}\mathbb{E}\left(\hf(\tc_n)\sum_{d(v)=\ell}\hf(\tc_{n,v})\right).
\]
However, we know that the height of $\st$ is $M$. Hence, for $\ell\geq M$, $\hf(\tc_n)$ is completely determined by $\tc_n^{(\ell)}$. Thus,
\[
\mathbb{E}\left(\hf(\tc_n)\sum_{d(v)=\ell}\hf(\tc_{n,v})\Big|\tc_n^{(\ell)}=T \wedge (\forall i \  (|T_i|=n_i) )\right)=\hf(T)\sum_{i=1}^{k}\E(\hf(\tc_{n_i}))=0.
\]
Therefore, we deduce using the law of total expectation that $\E\left(\hf(\tc_n)\Sigma_2\right)=0$. Putting all contributions together proves the lemma. 
\end{proof}

\end{document}
